\documentclass[aspectratio=169]{beamer}
\usepackage[T1]{fontenc}
\usepackage{graphicx}
\usepackage{xcolor}
\graphicspath{{figures/common/}{figures/04/}}
\definecolor{vaftblue}{RGB}{25,76,127}
\setbeamercolor{structure}{fg=vaftblue}
\setbeamerfont{title}{size=\fontsize{50}{56}\selectfont}
\setbeamerfont{subtitle}{size=\fontsize{24}{28}\selectfont}
\setbeamerfont{frametitle}{size=\fontsize{35}{40}\selectfont}
\setbeamerfont{normal text}{size=\fontsize{18}{22}\selectfont}
\setbeamertemplate{navigation symbols}{}
\setbeamertemplate{footline}[frame number]
\title{Fluctuations and Transient Events}
\subtitle{VAFT Introductory Tutorial | Session 04}
\author{}
\date{}

\begin{document}
\begin{frame}
  \titlepage
\end{frame}

\begin{frame}{Why this session exists}
  \begin{itemize}
    \item Session 03 handed you an equilibrium you can interrogate rather than assume.
    \item This one reads what will not hold still on top of it --- and says how much
      of that reading you are \emph{allowed} to claim.
  \end{itemize}
\end{frame}

\begin{frame}{A coil measures a derivative}
  \centering
  \[ S_{\mathrm{d}B/\mathrm{d}t}(f) = (2\pi f)^{2}\, S_{B}(f) \]
  \vspace{1em}

  An index fitted to raw pickup voltage is the field's \textbf{plus two}.
  VAFT never integrates behind an axis label.
\end{frame}

\begin{frame}{An array sets the answer, not the algorithm}
  \centering
  \[ \varphi_{\text{phase}} = \varphi_{0} - n\,\phi \]
  \vspace{0.5em}

  \begin{itemize}
    \item Two positions $120^{\circ}$ apart \(\rightarrow\) $n$ modulo 3.
    \item Three positions $90^{\circ}$ apart \(\rightarrow\) $n$ modulo 4.
  \end{itemize}
\end{frame}

\begin{frame}{Read the title before the number}
  \begin{itemize}
    \item The fit states how many \emph{distinct} toroidal positions answered,
      and which probes it left on another timebase.
    \item \texttt{available\_plots} says a call will not raise --- not that its
      answer is worth having.
  \end{itemize}
\end{frame}

\begin{frame}{Session workflow}
  \centering
  timing \(\rightarrow\) spectrogram \(\rightarrow\) spectrum and band power
  \(\rightarrow\) toroidal phase \(\rightarrow\) rational surfaces
\end{frame}

\begin{frame}{Where the chain stops}
  \begin{itemize}
    \item No resolver from $(m, n)$ to a surface; no rotation measurement, so no
      mode tracks; no event classifier of any kind.
    \item What you may write down: $f$ and its drift, the band power against a
      measured floor, $n$ modulo the array's alias step, $\psi_{N}$ of each bare
      $q$ surface --- and that the pairing between them is undetermined.
  \end{itemize}
\end{frame}
\end{document}
